\documentclass[11pt]{article}
\usepackage{amsmath,amssymb,amsthm,hyperref}
\newtheorem{proposition}{Proposition}
\newtheorem{lemma}{Lemma}
\title{Critical partition quantifiers: two invalid relaxations}
\author{The Eriksson Programme --- grouped research dossier}
\date{5 November 2026; unit 131; draft source, not a new paper}
\begin{document}
\maketitle
\paragraph{Status and scope.}
These scalar finite-torus arguments audit proposed shortcuts in the existing
rectangular-observation investigation. They do not refute the original
coefficient-adapted, coupled divergence-zero question. No new physical or
continuum conclusion, Lean verification, external assessment, numerical grade,
or originality claim is made. This source has not been rendered to PDF.

\paragraph{Conventions and the still-open question.}
Let $T_n=(\mathbb Z/n\mathbb Z)^D$, $D\ge3$, $n\ge3$, $N=n^D$,
$p=D/2$, and $q=2D/(D+2)$. All norms and sums use counting measure and
unit-spacing forward differences. Set
\[
 E_a(u)=\sum_{x\in T_n}\sum_{i=1}^D
 a(x)|u(x+e_i)-u(x)|^2.
\]
An allowed partition consists of products of nonempty cyclic intervals,
including whole cyclic coordinates. Its actual normalized mean map has
components $|B|^{-1/2}\sum_{x\in B}u(x)$. The orthogonal projection $P_{\cal P}$
onto block constants consequently satisfies
$\|P_{\cal P}u\|_2^2=\|C_{\cal P}u\|_2^2$.
The original order is: for each fixed admissible spatial $a$, construct
one partition depending on $a$ and then control every field. The hypotheses are
$1\le a\le2\mu$, $\operatorname{mean}a=\mu\ge1$ and
$\operatorname{mean}(\mu/a)^p\le M$; the proposed count is a constant times
$\max\{1,N/(1+\mu)^p\}$. Neither replacing the weighted class by a larger
unweighted critical ball nor choosing a partition before seeing $a$ is required
by that question.

\begin{proposition}[Critical-ball loss after forgetting the coefficient]
For every partition into $K$ nonempty cells (even nonrectangular),
\[
 \sup_{\substack{\sum_xu(x)=0\\\|\nabla u\|_q\le1}}
 \|(I-P_{\cal P})u\|_2
 \ \ge\ (2D)^{-1/q}\sqrt{1-K/N}.
\]
Thus uniformly vanishing errors on the full unweighted critical gradient ball
require $K/N\to1$. This is not a lower bound for the smaller weighted class
associated with a fixed coefficient.
\end{proposition}
\begin{proof}
For each vertex $x$ take $v_x=\delta_x-N^{-1}{\bf1}$. Its forward gradient
has exactly $2D$ nonzero occurrences, each of magnitude one; the restriction
$n\ge3$ avoids identifying the two endpoint neighbours. Hence
$\|\nabla v_x\|_q=(2D)^{1/q}$. Since constants lie in the projected space,
$(I-P_{\cal P})v_x=(I-P_{\cal P})\delta_x$. If $x\in B$,
\[
 \|(I-P_{\cal P})\delta_x\|_2^2=1-|B|^{-1},\qquad
 \sum_x\|(I-P_{\cal P})\delta_x\|_2^2=N-K.
\]
One of the $N$ terms is at least their average. Divide that witness by
$(2D)^{1/q}$. No compactness theorem or asymptotic continuum approximation
has been used.
\end{proof}

\paragraph{Why the weighted Holder step does not justify that relaxation.}
Holder on the directed edge occurrences, with exponents $2/q$ and
$2/(2-q)$, gives
\[
 \|\nabla u\|_q^2\le
 \left(D\sum_x a(x)^{-p}\right)^{1/p} E_a(u),
 \qquad \frac{q}{2-q}=p.
\]
It embeds one fixed weighted energy ball into a critical unweighted ball;
it does not identify the two. In the regime $\mu=\gamma n^2$, the upper
moment bound makes the factor at most $D^{1/p}M^{1/p}/\gamma$, which need
not vanish with $n$. For the constant profile $a=\gamma n^2$, however,
the exact Fourier Poincare inequality is
\[
 E_a(u)\ge4\gamma n^2\sin^2(\pi/n)\|u-\bar u\|_2^2
 \ge16\gamma\|u-\bar u\|_2^2.
\]
Here $\sin(\pi/n)\ge2/n$. For $\gamma\ge\max\{1,c/16\}$ the one-box
actual mean therefore already proves the desired scalar estimate with
coercivity $c$. The point witnesses in the enlarged ball do not contradict
this: their weighted energies grow like $n^2$.

\begin{proposition}[A coefficient-blind rectangular partition cannot suffice]
Fix $D\ge3$, $c>0$, and $\gamma\ge1$. Choose an integer $L\ge1$ with
\[
 \lambda_L=4D\sin^2\!\left(\frac{\pi}{2(L+1)}\right)<c,
 \qquad H=2L+3,
 \qquad m=2(L^D+D L^{D-1}).
\]
Such an $L$ exists, for instance whenever $D\pi^2/(L+1)^2<c$.
Set $M_0=m\gamma^p+1$. For every sufficiently large $n$, and every
rectangular partition selected before seeing the coefficient with $K<n/H$
cells, there is an admissible profile with
\[
 \mu=\gamma n^2,\qquad
 \operatorname{mean}(\mu/a)^p\le M_0,
 \qquad C_{\cal P}u=0,\qquad E_a(u)<c\|u\|_2^2
\]
for a nonzero scalar $u$. The same conclusion holds for any coarse penalty
$\Phi(C_{\cal P}u)$ with $\Phi(0)=0$.
In particular, there is no bounded-size partition depending only on
$(D,c,M_0,\mu,n)$ that works simultaneously for all profiles with these
parameters, even though $\max\{1,N/(1+\mu)^p\}=1$.
\end{proposition}
\begin{proof}
Some cell $B$ has volume at least $n^D/K$. Write its cyclic side lengths
as $\ell_i\le n$. For every $i$,
$\ell_i\ge |B|/n^{D-1}\ge n/K>H$.
Cut each interval at its own start (a whole cycle may be cut anywhere).
The cell contains a translated box $\{0,\ldots,2L+2\}^D$.

Inside it take $Q_1=\{1,\ldots,L\}^D$ and
$Q_2=Q_1+(L+2)e_1$. These cubes and their adjacent gradient-tail sets are
disjoint. On $Q_1$ define the positive Dirichlet product mode
\[
 v(j)=\prod_{i=1}^D\sin\left(\frac{\pi j_i}{L+1}\right),
 \qquad v=0\text{ off }Q_1.
\]
Take $u=v-\tau_{(L+2)e_1}v$. The equal translated sums cancel, and both
supports lie in $B$, so every block mean is zero. The elementary one-variable
identity $2s_j-s_{j-1}-s_{j+1}=4\sin^2(\pi/(2(L+1)))s_j$,
with $s_0=s_{L+1}=0$, gives the Dirichlet energy identity by multiplying
by $s_j$ and summing. Tensoring and adding the two disjoint copies gives
$E_1(u)=\lambda_L\|u\|_2^2$.

All nonzero forward differences have tails in
\[
 S=\bigcup_{r=1}^2\left(Q_r\cup\bigcup_{i=1}^D(Q_r-e_i)\right).
\]
For one cube this union consists of $L^D$ interior sites and $D$ disjoint
lower faces with $L^{D-1}$ sites each. Its size is $L^D+D L^{D-1}$;
the two unions are disjoint, so $|S|=m$. The padding prevents wrap
identifications; choose $n>H$ as part of ``sufficiently large.''
Set $f=m/N\le1/2$ and
\[
 a(x)=\begin{cases}1,&x\in S,\\ A,&x\notin S,\end{cases}
 \qquad A=\frac{\mu-f}{1-f},\quad \mu=\gamma n^2.
\]
Then $\operatorname{mean}a=\mu$ and $\mu\le A\le2\mu$, so the original
range holds. Since $2p=D$,
\[
 \operatorname{mean}(\mu/a)^p
 =f\mu^p+(1-f)(\mu/A)^p
 \le m\gamma^p+1=M_0.
\]
Every nonzero gradient tail is weak, including incoming boundary tails.
Consequently the full coupled energy, without discarding any edge, is
$E_a(u)=E_1(u)=\lambda_L\|u\|_2^2<c\|u\|_2^2$.
Finally $N/(1+\gamma n^2)^p<\gamma^{-p}\le1$.
\end{proof}

\paragraph{Quantifiers and limits of the conclusion.}
The last result proves $\forall{\cal P}\ \exists a\ \exists u$ for partitions
with $K<n/H$; it does not prove $\exists a\ \forall{\cal P}\ \exists u$.
The adversarial coefficient is explicitly chosen after the partition.
Allowing the partition to see $S$ escapes this argument. The witness is scalar;
no divergence-zero embedding is asserted. The coupled incoming weak edges are
included, unlike the distinct decoupled-energy obstruction of unit 129.
The result stops only the coefficient-blind shortcut, not the open route.

\paragraph{Primary-source applicability audit.}
Birman--Solomyak~\cite{BS67}, Theorem 2.1, uses a strictly positive
volume exponent for a lower-semiadditive cube functional. The discussion
immediately before Section 3 (printed page 302) explicitly notes failure at
zero exponent for a point load. Theorem 3.2 chooses partitions for each
function and assumes a strict subcritical target exponent. Theorems 3.3--3.4
instead provide measure-dependent partitions common to functions, with their
own strict hypotheses; the measure occurs in the target norm.
Thus neither a function-dependent approximation theorem nor that
measure-dependent theorem can simply be relabelled as our fixed spatial
gradient-weight problem. This applicability conclusion is our inference, not
a claim that no stronger result exists in the literature. In particular the
zero volume exponent is a known classical endpoint obstruction, not a new
discovery of this programme. The positive transverse exponent $D-k$ in unit
127 avoids that particular obstruction; invariance was still an extra
hypothesis there.

\paragraph{Model assessments and review context.}
This is a preparing-agent analytic self-audit in the current Codex session.
The system describes the assistant as GPT-6; the exact serving backend and
reasoning setting are not exposed by this session. No independent model,
human peer reviewer, external service, grade, executed diagnostic, or Lean
check is claimed. The checks consisted of deriving the trace identity,
Holder exponents, sine energy, full boundary-tail support, exact mean,
range, critical moment, and quantifier order. Originality remains
unestablished. No assessment of another artifact is inherited.

\paragraph{Next action and abandonment.}
Keep the single original mixed-cell/cross-boundary question open. A useful
construction must retain spatial information from the fixed $a$ and be
uniform over fields only after that coefficient is fixed. Before constructing
a manuscript, derive an actual overlap-to-disjoint-rectangle conversion with
both count and coupled-energy control, or document the concrete obstruction.
Do not replace this obligation by an assumed partition certificate. This
audit does not resolve that obligation; conservatively count it as the second
unit without advance on the original question, following verification unit
130. At the third such unit, stop/change the mechanism or record its precise
blockage. Group these negative guards with the existing dossier; do not create
a separate paper, deposit, or regrade the published v5.

\begin{thebibliography}{1}
\bibitem{BS67} M. S. Birman and M. Z. Solomyak,
\emph{Piecewise-polynomial approximations of functions of the classes
$W_p^\alpha$}, Math. USSR-Sbornik 2 (1967), 295--317.
\href{https://doi.org/10.1070/SM1967v002n03ABEH002343}{doi:10.1070/SM1967v002n03ABEH002343}.
\href{https://www.mathnet.ru/php/getFT.phtml?jrnid=sm\&paperid=4132\&what=fullteng}{Primary full text}.
\end{thebibliography}
\end{document}
